\section{工程实践与调参}

\subsection{Replay buffer 和数据采样策略}

现代 off-policy Actor-Critic 模型往往使用优先级 replay buffer、frame stacking 和 multi-environment rollout。过往讲义中也反复强调：数据采样策略直接决定了 critic 的训练信号。

\begin{knowledgebox}{Prioritized replay 的小技巧}
把 TD error 的绝对值作为 priority，在 sampling 时增加概率，并在 loss 中乘以 importance weight correction，确保 update 不会偏向高 error 样本。
\end{knowledgebox}

\subsection{Batching 与多步回报的实践}

批量大小、采样频率与 n-step 的组合是基础调参方向。可以通过下表把常用组合总结出来。

\begin{table}[H]
\centering
\begin{tabular}{p{0.25\textwidth}p{0.35\textwidth}p{0.28\textwidth}}
\toprule
配置项 & 常见取值 & 训练影响 \\
\midrule
Batch size & 256--4096 transitions & 越大 variance 越低，但硬件吞吐是瓶颈 \\
n-step & 1--5 & 大步数偏差小但有 drift，常用 $\lambda$ mix \\
Update frequency & 1 update per 1--4 steps & 更少 update 减少 off-policy bias \\
Replay ratio & 1--4 & >1 表示每次采样后更新多次，各步 learning rate 需对应减小 \\
\bottomrule
\end{tabular}
\caption{实践中常调的超参组合}
\end{table}

\subsection{调参与故障排查}

\begin{figure}[H]
\centering
\includegraphics[width=0.82\textwidth]{slide-15.jpg}
\caption{Lecture 4 中展示的训练流程图：多 worker 拉取数据、critic 拟合目标、Actor 更新策略。}
\end{figure}

当训练出现 divergence，先看 three signals：critic loss、KL、reward curve。此外，常见的failure mode 有：
\begin{itemize}
    \item Critic overestimate：回归失灵导致 advantage 恒正；
    \item Replay buffer stale：环境改变但 buffer 中数据仍旧；
    \item Exploration collapse：entropy 快速降为 0。
\end{itemize}

\begin{warningbox}{出现 divergence 时的快速复位流程}
先把策略退回上一个 checkpoint，关闭 entropy decay，减小 learning rate，再逐步恢复。不要一口气加大数据量，否则会加速灾难。
\end{warningbox}

\subsection{本章小结}

实践中，数据采样、batching 和超参调度构成 Actor-Critic 稳定训练的系统性问题。监控 replay ratio、entropy 和 critic loss 是调参的第一手数据。

